\documentclass[10pt]{article}

\usepackage[utf8]{inputenc}

\usepackage[T1]{fontenc}

\usepackage{amsmath}

\usepackage{amsfonts}

\usepackage{amssymb}

\usepackage[version=4]{mhchem}

\usepackage{stmaryrd}



\begin{document}

Собств нн о е значение к в а д р а тной м а т р ц ц ы $A$ над полем $k$ (или характеристический корень) - корень ее характеристического многочлена.



Лит. см. при статьях Линейное преобразование, Матрича. СОБСТВЕННЫЕ ЗНАЧЕНИЯ ДИФФЕРЕНЦИАЛЬНЫХ ОПЕРАТОРОВ; ч и с л е н ны е ме т о ды на$\mathrm{x} о$ ж д е н и я - методы вычисления собственных значений и соответствующих собственных функций дифференциальных операторов. Колебания упругих ограниченных тел ошисываются уравнением



\[

\frac{\partial^{2} \varphi}{\partial t^{2}}=L \varphi

\]



где $L \varphi$ - нек-рое диффференциальное выражение. Если решение уравнения (1) искать в виде



\[

\varphi=T(t) u(x)

\]



то относительно функции $u$ получается уравнение



\[

L(u)+\lambda u=0

\]



внутри ограниченной области при нек-рых однородных условиях на ее границе. Значения параметра $\lambda$, при к-рых существуют отличные от тождественного нуля решения уравнсния (2), удовлетворяющие однородным краевым условиям, наз. с о б с т в е н н ы м и $\mathbf{3}$ н ач е ния м и (чис л а ми), а их соответствующие репения - с об с т в е н ным и ф у н к ци я м и. Возникающая при этом диффференциальная задача на собственные значения состоит в нахождении собственных значений $\lambda$ и соответствующих им собственных функций.



Численное решение диффференциальной задачи на собственные значения проводится в три этапа:



\begin{enumerate}

  \item сведение задачи к более простой, напр. к алгебраической (дискретной);



  \item выяснение точности дискретной задачи;



  \item вычисление собственных значений дискретной задачи (см. Линейнал алгебра; численные методы).



\end{enumerate}



Сведение к дискретной задаче. Сведение задачи (2) к ее дискретной модели производится в основном сеток методом и проекционными методами. При этом естественно требовать, чтобы основные свойства исходной задачи сохранялись в ее дискретном аналоге. В частности, должна сохраняться самосопряженность соответствующих дискретных огераторов в пространстве фуункций дискретного аргумента.



Одним из методов такого сведения является и н т ег р о-и н т е р по ля ци онны й ме тод. Напр., пусть поставлена задата



\[

\begin{gathered}

\frac{d}{d x}\left(\frac{1}{p(x)} \frac{d u}{d x}\right)+\lambda u=0, \quad 0<x<1, \\

u(0)=0, \quad u(1)=0 .

\end{gathered}

\]



Эта задача возникает, напр., при изучении поперечных колебаниї неоднородной струны и продольных колебаний неоднородного стержня.



На отрезке $[0,1]$ вводится разностная сетка $\bar{\omega} \mathrm{c}$ узлами $x_{i}=i h, i=0,1, \ldots, N, h=1 / N$. Каждому узлу $x_{i}, i=1, \ldots, N-1$, ставится в соответствие әлементарная область $S_{i}\left\{x_{i}-h / 2 \leqslant x \leqslant x_{i}+h / 2\right\}$. Интегрирование по областям $S_{i}$ уравнения (3) приводит к выражению



\[

\left.\begin{array}{c}

\frac{w_{i-1 / 2}-w_{i+1 / 2}}{h}=\lambda \frac{1}{h} \int_{x_{i}-h / 2}^{x_{i}+h / 2} u d x, \\

w_{i \pm 1 / 2}=\left(\frac{1}{p} \frac{d u}{d x}\right)_{x=x_{i} \pm h / 2} .

\end{array}\right\}

\]



Пусть



\[

\begin{gathered}

u=\mathrm{const}=u_{i}, \quad x_{i}-h / 2 \leqslant x \leqslant x_{i}+h / 2, \\

w=\mathrm{const}=w_{i-1 / 2}, \quad x_{i-1} \leqslant x \leqslant x_{i} .

\end{gathered}

\]



Тогда



\[

w_{i-1 / 2}=\frac{u_{i}-u_{i-1}}{h a_{i}}, \quad a_{i}=\frac{1}{h} \int_{x_{i-1}}^{x_{i}} p(x) d x .

\]



При подстановке (6) в (5) получается уравнение $-\frac{1}{h}\left(\frac{v_{i+1}-v_{i}}{a_{i+1} h}-\frac{v_{i}-v_{i-1}}{a_{i} h}\right)=\lambda h_{v_{i}}, \quad i=1,2, \ldots, N-1$,



где $v$ - искомая сеточная функция. Краевые условия $v_{0}=0, v_{N}=0$ приводят к алгебраич. задаче на собственные значения



\[

A v=\lambda^{h} v

\]



где $A$ - трехдиагональная симметрич. матрица порядка $n=N-1$ (см. [10]).



В а р и а І и о нн 0-р а 3 ност ны й ме тод сведения к дискретной задаче используется, когда задача на собственные значения моякет быть сформулирована как вариационная. Напр., собственные значения задачи (3), (4) являются стационарными значениями функционала



\[

\frac{D[u]}{H[u]}, \quad D[u]=\int_{0}^{1} \frac{1}{p(x)}\left(\frac{d u}{d x}\right)^{2} d x ; \quad H[u]=\int_{0}^{1} u^{2} d x

\]



При замене интегралов квадратурными суммами, а производных - разностными отношениями, дискретныї аналог функционала имеет вид



\[

\begin{gathered}

\frac{D^{h}[v]}{H^{h}[v]}, D^{h}[v]=\sum_{i=1}^{N} \frac{1}{a_{i}}\left(\frac{v_{i}-v_{i-1}}{h}\right)^{2} h ; \\

H^{h}[v]=\sum_{i=0}^{N} v_{i}^{2} h,

\end{gathered}

\]



где $a_{i}$ - разностный аналог коэффициента $p(x)$, к-рый может быть вычислен по формуле (6). Дискретный аналог задачи (3), (4) получается из необходимого условия әкстремума



\[

\frac{\partial}{\partial v_{i}}\left(D^{h}[v]-\lambda^{h} H^{h}[v]\right)=0, \quad i=1,2, \ldots, N-1 .

\]



Дифференцирование ігриводит снова к задаче (7) (см. [9]).



Проекционно-ра сведения $\kappa$ дискретной задаче состоит в следующем. Выбирается линейно независимая координатная система функций $\alpha_{i}, i=1,2, \ldots, n$, и линейно пезависимая проекционная система функций $\beta_{j}, j=1,2, \ldots, n$. Приближкеные собственные функций ищутся в виде



\[

\tilde{u}=\sum_{i=1}^{n} v_{i} \alpha_{i}

\]



Коэффициенты разложения $v_{i}$ и приближенныс собственпые значения определяются из условия



\[

\left(L \tilde{u}+\tilde{\lambda} \tilde{u}, \beta_{j}\right)=0, \quad j=1,2, \ldots, n,

\]



где (, ) - скалярное произведение в гильбертовом пространстве. [При совпадении координатной и проекционной систем говорят о методе Бубнова - Галеркина. Если, кроме того, оператор дифференциальной задачи самосопряженный, то метод наз. методом Рэлея Ритца (см. [4]). ] В частности, для задачи (3), (4), если все $\alpha_{i}=\beta_{i}$ удовлетворяют (4), условие (8) принимает вид



\[

\begin{gathered}

\sum_{i=1}^{n} v_{i} \int_{0}^{1}\left(\frac{1}{p(x)} \frac{d \alpha_{i}}{d x} \frac{d \alpha_{j}}{d x}-\tilde{\lambda} \alpha_{i} \alpha_{j}\right) d x=0, \\

j=1, \ldots, n .

\end{gathered}

\]



Чтобы упростить получение алгебраич. задачи, систему функций $\alpha_{i}$ выбирают почти ортогональной.





\end{document}